% ============================================================
% The Centered Collision Sum
% Combined research note (merging collision transform and
% centered collision sum papers)
% ============================================================

\documentclass[12pt]{amsart}

\usepackage{amsmath,amssymb,amsthm}
\usepackage{booktabs}
\usepackage{microtype}
\usepackage{url}

\newtheorem{theorem}{Theorem}
\newtheorem{lemma}[theorem]{Lemma}
\newtheorem{proposition}[theorem]{Proposition}
\newtheorem{corollary}[theorem]{Corollary}
\newtheorem{conjecture}[theorem]{Conjecture}
\newtheorem{observation}[theorem]{Observation}
\theoremstyle{definition}
\newtheorem{definition}[theorem]{Definition}
\theoremstyle{remark}
\newtheorem{remark}[theorem]{Remark}

\title{The Centered Collision Sum}

\author{Alexander S.\ Petty}

\email{alexander.petty@gmail.com}
\date{October 2023; revised April 2026}
\subjclass[2020]{11A63, 11N05, 11M06}

\begin{document}
\begin{abstract}
For a prime $p > b$ and lag $\ell \ge 1$, the collision
deviation $S_{\ell}(p) = C(b^{\ell} \bmod p) - Q$ depends
only on $p \bmod b^{\ell+1}$ (the finite determination
theorem). This reduces the collision deviation to a function
on the finite group
$(\mathbb{Z}/b^{\ell+1}\mathbb{Z})^{\times}$.

The \emph{collision transform} decomposes this function over
Dirichlet characters modulo~$b^{\ell+1}$. The trivial
coefficient is exactly $-1/2$, forced by the reflection
identity $S(a) + S(m{-}a) = -1$. The non-trivial characters
carry the centered fluctuation.

Subtracting the per-class mean defines the centered
fluctuation sum $F^{\circ}(1, \ell) = \sum_p S^{\circ}_p / p$.
By Mertens' theorem for arithmetic progressions, $F^{\circ}$
converges at $s = 1$ for every base~$b$ and lag~$\ell$. This
is the Centered Collision PNT: the centered collision
fluctuations cancel at the edge of the critical strip.

For prime bases at lag~$1$, finite evaluation suggests the
explicit class mean formula
$\overline{S}_R = (R{+}1)/b - 1$, compatible with grand mean
$-1/2$ and reflection symmetry
$\overline{S}_R + \overline{S}_{b-2-R} = -1$.
\end{abstract}
\maketitle

% ============================================================
\section{Introduction}

This paper proves the convergence of the centered collision
fluctuation sum. The preceding paper~\cite{paper-fluct}
conjectured that the raw fluctuation sum $F(1, \ell)$ diverges
at the Mertens rate $-\mu_{\ell} \cdot \log\log x$, driven by
a persistent per-class bias. The present paper identifies the
bias exactly, subtracts it, and proves that the remainder
converges.

The proof has three steps. First, the finite determination
theorem (Section~2) shows that $S_{\ell}(p)$ depends only on
$p \bmod b^{\ell+1}$. Second, the collision transform
(Section~3) decomposes the resulting finite function over
Dirichlet characters. Third, the Centered Collision PNT
(Section~6) assembles the convergence from Mertens' theorem
for arithmetic progressions.

% ============================================================
\section{The Finite Determination Theorem}

\begin{theorem}[Finite determination]\label{thm:finite}
For any base $b \ge 2$, lag $\ell \ge 1$, and integer
$p > b^{\ell+1}$ with $\gcd(p, b) = 1$, the collision
deviation $S_{\ell}(p) = C(b^{\ell} \bmod p) - \lfloor
(p{-}1)/b \rfloor$ depends only on $p \bmod b^{\ell+1}$.
Primality of $p$ is not required.
\end{theorem}

\begin{proof}
Set $m = b^{\ell+1}$. For each $r \in \{1, \ldots, p{-}1\}$,
define the slice index $n(r) = \lfloor mr/p \rfloor
\in \{0, \ldots, m{-}1\}$.

Since the slice width $p/m > 1$, each slice contains integers.
The digit function and the collision multiplier are constant on
each slice:
\[
\delta(r) = \left\lfloor \frac{br}{p} \right\rfloor
= \left\lfloor \frac{n}{b^{\ell}} \right\rfloor,
\qquad
\delta(b^{\ell} r \bmod p) = n \bmod b.
\]

A slice is \emph{diagonal} if
$\lfloor n/b^{\ell} \rfloor = n \bmod b$. Let
$G_{\ell,b} = \{n \in \{0, \ldots, m{-}1\} :
\lfloor n/b^{\ell} \rfloor = n \bmod b\}$.
Then $|G_{\ell,b}| = b^{\ell}$ and
\[
C(b^{\ell} \bmod p) = -1 + \sum_{n \in G_{\ell,b}}
\left(
\left\lfloor \frac{(n{+}1)p}{m} \right\rfloor
- \left\lfloor \frac{np}{m} \right\rfloor
\right).
\]

Write $p = mt + a$ with $1 \le a < m$ and $\gcd(a, b) = 1$.
Then
\[
\left\lfloor \frac{(n{+}1)p}{m} \right\rfloor
- \left\lfloor \frac{np}{m} \right\rfloor
= t + \left\lfloor \frac{(n{+}1)a}{m} \right\rfloor
- \left\lfloor \frac{na}{m} \right\rfloor.
\]
Since $Q = b^{\ell} t + \lfloor a/b \rfloor$:
\[
S_{\ell}(p) = -1 - \left\lfloor \frac{a}{b} \right\rfloor
+ \sum_{n \in G_{\ell,b}}
\left(
\left\lfloor \frac{(n{+}1)a}{m} \right\rfloor
- \left\lfloor \frac{na}{m} \right\rfloor
\right),
\]
which depends only on $a = p \bmod m$.
\end{proof}

\begin{remark}
The modulus $b^{\ell+1}$ is sharp. In base~$10$ at lag~$2$:
$S_2(1009) = 8$ and $S_2(1109) = -10$, though
$1009 \equiv 1109 \pmod{100}$. Reduction modulo~$b^2$ does
not determine $S_2$; modulo $b^3 = 1000$ does.
\end{remark}

% ============================================================
\section{The Collision Transform}

Since $S_{\ell}$ is a function on
$(\mathbb{Z}/b^{\ell+1}\mathbb{Z})^{\times}$, it admits a
Fourier expansion over Dirichlet characters.

\begin{definition}
The \emph{collision transform} is
\[
\hat{S}_{\ell}(\chi) = \frac{1}{\phi(b^{\ell+1})}
\sum_{a \in (\mathbb{Z}/b^{\ell+1}\mathbb{Z})^{\times}}
S_{\ell}(a)\, \overline{\chi}(a)
\]
for each Dirichlet character $\chi$ modulo~$b^{\ell+1}$.
\end{definition}

\begin{theorem}[Reflection identity]\label{thm:reflection}
For $m = b^{\ell+1}$ and any unit $a$:
\[
S_{\ell}(a) + S_{\ell}(m - a) = -1.
\]
\end{theorem}

\begin{proof}
Replacing $a$ by $m - a$ in the formula of
Theorem~\ref{thm:finite} sends
$\lfloor na/m \rfloor$ to $n - 1 - \lfloor na/m \rfloor$
for $0 < na < m$ (the complementary floor identity). The
sum over diagonal slices transforms accordingly, giving
$S_{\ell}(m - a) = -1 - S_{\ell}(a)$.
\end{proof}

\begin{corollary}[Trivial coefficient]\label{cor:trivial}
$\hat{S}_{\ell}(\chi_0) = -1/2$.
\end{corollary}

\begin{proof}
The involution $a \mapsto m - a$ pairs the units. By
Theorem~\ref{thm:reflection}, each pair averages to $-1/2$.
\end{proof}

\begin{corollary}[Nontrivial even characters vanish]
For every nontrivial even character $\chi$ (i.e., $\chi(-1) = +1$):
$\hat{S}_{\ell}(\chi) = 0$.
\end{corollary}

\begin{proof}
Since $\chi$ is even, $\chi(-a)=\chi(a)$. Pairing $a$ with
$-a$ and using Theorem~\ref{thm:reflection},
\[
\hat{S}_{\ell}(\chi)
= \frac{1}{2\phi(b^{\ell+1})}
\sum_{a}
\bigl(S_{\ell}(a)+S_{\ell}(-a)\bigr)\overline{\chi}(a)
= -\frac{1}{2\phi(b^{\ell+1})}\sum_a \overline{\chi}(a).
\]
For a nontrivial character, the last sum is zero.
\end{proof}

% ============================================================
\section{The Half-Group Structure}

\begin{theorem}[Half-group]\label{thm:half}
For every diagonal slice $n \in G_{\ell,b}$ with $n \neq m-1$,
the wrapping set
$W_n = \{a : (n{+}1)a \bmod m < a\}$
satisfies $|W_n| = \phi(m)/2$.
\end{theorem}

\begin{proof}
Let $c = n + 1$. If $a \in W_n$, then
$c(m{-}a) \bmod m = m - (ca \bmod m) > m - a$,
so $m - a \notin W_n$. The involution $a \mapsto m - a$
maps $W_n$ bijectively onto its complement.
\end{proof}

% ============================================================
\section{The Class Mean Formula}

\begin{definition}
For spectral class $R = (p{-}1) \bmod b$, the class mean is
\[
\overline{S}_R = \frac{1}{|\mathcal{A}_R|}
\sum_{a \in \mathcal{A}_R} S_{\ell}(a),
\]
where $\mathcal{A}_R = \{a \in
(\mathbb{Z}/b^{\ell+1}\mathbb{Z})^{\times} :
a \equiv R{+}1 \pmod{b}\}$.
\end{definition}

\begin{observation}[Prime-base class mean at lag~$1$]
\label{prop:class-mean}
For prime base $b$ at lag $\ell = 1$:
$\overline{S}_R = (R{+}1)/b - 1$.
\end{observation}

This is obtained by evaluating $S_1(a)$ at all
$a \in (\mathbb{Z}/b^2\mathbb{Z})^{\times}$ using
Theorem~\ref{thm:finite} and averaging over each class.

For composite bases, the class means are rational and
satisfy the following universal properties.

\begin{corollary}\label{cor:universal}
For all bases $b \ge 2$ at any lag $\ell \ge 1$, let
$\mathcal{R} = \{R \in \{0, \ldots, b{-}2\} : \gcd(R{+}1, b) = 1\}$
be the set of admissible spectral classes. Then:
\begin{enumerate}
\item[\textup{(i)}] Grand mean:
$\displaystyle\frac{1}{|\mathcal{R}|}
\sum_{R \in \mathcal{R}} \overline{S}_R = -\frac{1}{2}$.
\item[\textup{(ii)}] Reflection symmetry:
$\overline{S}_R + \overline{S}_{b-2-R} = -1$
for every $R \in \mathcal{R}$.
\end{enumerate}
\end{corollary}

\begin{proof}
Property~(i) follows from
Corollary~\ref{cor:trivial}. Property~(ii)
follows from the reflection identity: the map
$a \mapsto m - a$ sends class~$R$ to class~$b{-}2{-}R$.
\end{proof}

\begin{table}[h]
\centering
\begin{tabular}{rrrrr}
\toprule
$R$ & $0$ & $2$ & $6$ & $8$ \\
\midrule
$\overline{S}_R$ (exact) & $-17/10$ & $-9/10$ & $-1/10$ & $7/10$ \\
computed & $-1.709$ & $-0.897$ & $-0.097$ & $+0.699$ \\
\bottomrule
\end{tabular}
\caption{Class means in base~$10$ at lag~$1$, computed over
$166{,}000$ primes per class ($664{,}574$ total, $p \le 10^7$).}
\label{tab:class-means}
\end{table}

% ============================================================
\section{The Centered Collision PNT}

\begin{definition}
The centered collision deviation is
$S^{\circ}_p = S_{\ell}(p) - \overline{S}_{R(p)}$.
The centered fluctuation sum is
\[
F^{\circ}(s, \ell) = \sum_{\substack{p > b^{\ell+1} \\
\gcd(p,b) = 1}} \frac{S^{\circ}_p}{p^s}.
\]
\end{definition}

\begin{theorem}[Centered Collision PNT]\label{thm:cpnt}
For every base $b \ge 2$ and lag $\ell \ge 1$,
$F^{\circ}(1, \ell)$ converges.
\end{theorem}

\begin{proof}
Set $q=b^{\ell+1}$. By Theorem~\ref{thm:finite}, $S_{\ell}(p)$
depends only on the residue class of $p$ modulo $q$, so the
same is true for the centered quantity $S_p^\circ$. Define the
finite function
\[
f^\circ(a)=S_\ell(a)-\overline{S}_{R(a)}
\qquad
\bigl(a\in(\mathbb Z/q\mathbb Z)^\times\bigr),
\]
where $R(a)$ is determined by $a\equiv R(a)+1 \pmod b$.
By definition of the class means,
\[
\sum_{a\in \mathcal A_R} f^\circ(a)=0
\]
for every spectral class $\mathcal A_R$, hence in particular
\[
\sum_{a\in (\mathbb Z/q\mathbb Z)^\times} f^\circ(a)=0.
\]
Therefore the trivial Fourier coefficient vanishes, and the
collision transform expands as
\[
f^\circ(a)=\sum_{\chi \ne \chi_0}\hat{f}^{\circ}(\chi)\,\chi(a),
\]
with a finite sum over nontrivial Dirichlet characters modulo $q$.
Substituting $a=p \bmod q$ gives
\[
F^{\circ}(1, \ell)
= \sum_{\chi \ne \chi_0}\hat{f}^{\circ}(\chi)
\sum_{\substack{p>b^{\ell+1}\\ p\nmid b}} \frac{\chi(p)}{p}.
\]

For each nontrivial $\chi$ modulo $q$, Mertens' theorem for
arithmetic progressions (see, for example, the discussion in
\cite[Ch.~22]{davenport}) gives
\[
\sum_{\substack{p \le x \\ p \equiv a \pmod{q}}}
\frac{1}{p} = \frac{\log\log x}{\phi(q)} + M_a + o(1).
\]
Summing this identity against $\chi(a)$ over the reduced residue
classes modulo $q$ yields
\[
\sum_{p \le x} \frac{\chi(p)}{p}
\;=\;
\frac{\log\log x}{\phi(q)}\sum_a \chi(a)
+ \sum_a \chi(a) M_a + o(1).
\]
Because $\chi$ is nontrivial, $\sum_a \chi(a)=0$, so the
$\log\log x$ term cancels. Hence each prime sum
\[
\sum_{p} \frac{\chi(p)}{p}
\]
converges. Since only finitely many characters occur, the outer
sum is finite, and therefore $F^\circ(1,\ell)$ converges.
\end{proof}

\begin{table}[h]
\centering
\begin{tabular}{rrrrrr}
\toprule
primes & $p_{\max}$ & $F(1,1)$ & $F(1,1)/\!\log\!\log p$
& $F^{\circ}$ & $F^{\circ}/\!\log\!\log p$ \\
\midrule
$1{,}000$ & $7{,}949$ & $-1.46$ & $-0.664$
& $0.020$ & $0.009$ \\
$10{,}000$ & $104{,}773$ & $-1.69$ & $-0.690$
& $0.023$ & $0.009$ \\
$100{,}000$ & $1{,}299{,}791$ & $-1.86$ & $-0.705$
& $0.023$ & $0.009$ \\
$664{,}574$ & $9{,}999{,}991$ & $-1.99$ & $-0.714$
& $0.023$ & $0.008$ \\
\bottomrule
\end{tabular}
\caption{Raw vs.\ centered fluctuation sums in base~$10$
at lag~$1$. The raw ratio drifts toward a Mertens constant.
The centered ratio decreases toward zero, consistent with
convergence.}
\label{tab:centered}
\end{table}

% ============================================================
\section{Computational Universality Checks}

\begin{table}[h]
\centering
\begin{tabular}{rrrrrr}
\toprule
base & lag & primes & $F(1,\ell)$ & $F^{\circ}$ & behavior \\
\midrule
$3$ & $1$ & $18{,}000$ & $-1.59$ & $0.163$ & small \\
$7$ & $1$ & $42{,}000$ & $-1.66$ & $0.031$ & small \\
$10$ & $1$ & $78{,}000$ & $-1.76$ & $0.043$ & small \\
$10$ & $2$ & $78{,}000$ & $-1.83$ & $0.072$ & small \\
$10$ & $3$ & $78{,}000$ & $-0.56$ & $-0.227$ & small \\
$12$ & $1$ & $42{,}000$ & $-1.67$ & $0.129$ & small \\
\bottomrule
\end{tabular}
\caption{Centered fluctuation sums across bases and lags.
All centered sums remain comparatively small over the computed
range.}
\label{tab:universal}
\end{table}

% ============================================================
\section{Computational Resonance Data}

Define $\psi_k(p) = e(kc_p/p)$, where
$c_p = b(1 - b^{\ell})^{-1} \bmod p$ is the collision
parameter. The partial sums
$\Sigma_k(x) = \sum_{p \le x} \psi_k(p)$ exhibit a
dichotomy controlled by $\gcd(k, b{-}1)$.

\begin{table}[h]
\centering
\begin{tabular}{rrrr}
\toprule
$k$ & $\gcd(k, b{-}1)$ & $|\Sigma_k| / \sqrt{\pi(x)}$
& behavior \\
\midrule
$1$ & $1$ & $0.16$ & cancels \\
$2$ & $1$ & $0.21$ & cancels \\
$3$ & $3$ & $140$ & resonates \\
$5$ & $1$ & $0.06$ & cancels \\
$7$ & $1$ & $0.22$ & cancels \\
$9$ & $9$ & $280$ & resonates \\
$10$ & $1$ & $0.15$ & cancels \\
\bottomrule
\end{tabular}
\caption{Fourier modes in base~$10$ at $78{,}493$ primes.
Modes with $\gcd(k, b{-}1) = 1$ are numerically small at this
scale, while modes with $\gcd(k, b{-}1) > 1$ are numerically
large.}
\label{tab:resonance}
\end{table}

The data suggest that the resonant modes correspond to
characters modulo~$b^{\ell+1}$ that factor through characters
modulo~$b$, while the non-resonant modes distinguish
sub-classes within each spectral class.

% ============================================================
\section{Remarks}

\subsection*{Connection to the alignment limit}

The class mean $\overline{S}_0 = -2/3$ in base~$3$ is the
alignment limit of the prime~$3$ from~\cite{paper1}. The
constant $2/3$ appears both as the alignment limit and as
the magnitude of the class bias in the simplest nontrivial
case. More generally, for prime base~$b$,
$\overline{S}_0 = 1/b - 1 = -(b{-}1)/b$.

\subsection*{The Mertens constant}

The grand mean $-1/2$ controls the leading-order Mertens
constant. By Mertens' theorem for arithmetic progressions,
the raw sum diverges at rate $(-1/2) \cdot \log\log x$
from the grand mean alone. The per-class deviations from
$-1/2$ and the spectral class corrections determine the
exact Mertens constant $\mu_{\ell}$.

\subsection*{Toward an explicit collision PNT}

A genuine collision analogue of the prime number theorem
would give an asymptotic formula for the unweighted sum
$\sum_{p \le x} S^{\circ}(p)$ with a proved error bound.
The collision transform expresses $S^{\circ}$ as a linear
combination of Dirichlet characters, and the classical
explicit formula for character sums provides the machinery.

% ============================================================
\begin{thebibliography}{99}

\bibitem{paper1}
A.~S.~Petty,
Why the golden ratio selects the prime three,
research note, 2020.

\bibitem{paper-fluct}
A.~S.~Petty,
The collision fluctuation sum,
research note, 2023.

\bibitem{hardy-wright}
G.~H.~Hardy and E.~M.~Wright,
\emph{An Introduction to the Theory of Numbers},
6th~ed., Oxford University Press, 2008.

\bibitem{mauduit-rivat}
C.~Mauduit and J.~Rivat,
Sur un probl\`eme de Gelfond: la somme des chiffres des
nombres premiers,
\emph{Ann.\ of Math.}\ \textbf{171} (2010), 1591--1646.

\bibitem{davenport}
H.~Davenport,
\emph{Multiplicative Number Theory},
3rd~ed., Springer, 2000.

\bibitem{mertens}
F.~Mertens,
Ein Beitrag zur analytischen Zahlentheorie,
\emph{J.\ Reine Angew.\ Math.}\ \textbf{78} (1874), 46--62.

\bibitem{shparlinski}
I.~E.~Shparlinski,
Character sums over shifted primes,
\emph{Mathematika} \textbf{56} (2010), 363--372.

\bibitem{sos}
V.~T.~S\'os,
On the distribution mod~1 of the sequence $n\alpha$,
\emph{Ann.\ Univ.\ Sci.\ Budapest.\ E\"otv\"os Sect.\ Math.}
\textbf{1} (1958), 127--134.

\bibitem{stanley}
R.~P.~Stanley,
\emph{Enumerative Combinatorics}, vol.~1,
2nd~ed., Cambridge University Press, 2012.

\bibitem{nfield}
A.~S.~Petty,
\emph{nfield: A structural analysis engine for fractional
field invariants},
software, 2025--2026,
tag \texttt{boundary-alphabet-q2-v4-2026-04-10-r2},
commit \texttt{e0b799d8ab947f09b2d714dfe2a70ea20f7251e3}.
Repository: \url{https://github.com/alexspetty/nfield}.

\end{thebibliography}

\end{document}
